% Relocated narrative from ch06_reliability.tex. Chapter aliases live on ch_productization.tex. Do not re-add \chapter or ch:* labels here.


\textit{Read this reference section for:} what a qualification claim is, mechanism-driven
planning, acceleration validity, sample confidence, acceptance, and the bounded
decisions qualification supports.

\textit{Use the technical chapters for:} source physics, wavelength control,
packaging, and receiver mechanisms
(\Cref{ch:lasers,ch:wdm,sec:laser-aging,sec:pd-tia}).

\textit{Use the reliability reference appendix for:} standards detail, named
stress methods, connector test methods, and the planning matrix
(\Cref{app:reliability-reference}).

A five-year lifetime requirement cannot be verified by operating the product
for five years before release. Reliability engineering therefore works
indirectly. You identify the physical mechanisms that could cause the product
to violate its claim, select stresses that accelerate those mechanisms, measure
signatures that reveal change, and decide whether the resulting evidence
supports the intended use.

That chain is the whole discipline: claim, mechanism, stress, observable,
acceptance, samples, confidence, decision. Every link has to hold. A stress
without a named mechanism is only exposure. An observable chosen after the
stress is a story told backwards. A sample that covers one lot bounds one lot.
Running more tests does not repair a broken link, and omitting a test is not
automatically irresponsible if no credible mechanism sits behind it.

Two limits follow, and both matter more than any individual test result. First,
qualification bounds confidence for the tested design, sample population,
mechanisms, stresses, durations, observables, and assumptions. It does not
prove that every future production unit will be reliable.
\Cref{ch:manufacturing} asks whether production can reproduce that qualified
population consistently. Second, the fleet is what actually tests the claim.
Qualification buys you the right to ship with a stated residual risk, and the
field either confirms the mechanism model or sends back evidence that it was
wrong (\Cref{ch:failure-modes}).

\section{What reliability qualification proves}
\label{sec:qual-terms}

Several activities use the same chambers, the same instruments, and sometimes
the same units. They are not the same job. What separates them is the question
asked, the population exposed, the observable recorded, and the decision the
result unlocks.

\begin{table*}[ht]
\refstepcounter{table}\label{tab:qual-activity-split}%
\centering
\setlength{\tabcolsep}{3pt}%
\footnotesize
\renewcommand{\arraystretch}{1.25}%
\begin{tabular}{@{}p{0.22\linewidth}p{0.40\linewidth}p{0.34\linewidth}@{}}
\toprule
Activity & Question answered & Main output \\
\midrule
Characterization & How does a healthy product behave across corners and units? & Behavior maps and distributions \\
\addlinespace[1.0em]
System validation & Is the complete product suitable for its intended system use? & Suitability and margin evidence across the supported envelope \\
\addlinespace[1.0em]
Reliability qualification & Do named mechanisms violate the life and environmental claim? & Bounded confidence statement plus residual risk \\
\addlinespace[1.0em]
Manufacturing validation & Can production reproduce and control the qualified design? & Process capability, yield, and control evidence \\
\addlinespace[1.0em]
Production screening & Does this unit, or this sampled output, show a detectable defect? & Accept, reject, or contain decisions per unit or lot \\
\bottomrule
\end{tabular}
\par\medskip
{\raggedright\footnotesize\textbf{Table~\thetable.} The same temperature
chamber can support several of these activities. The category is set by the
claim, the exposure, the population, the observable, and the decision, not by
the equipment. Lifecycle ownership: \Cref{ch:product-readiness,ch:manufacturing}.\par}
\end{table*}
\FloatBarrier

The distinction that gets missed most often is characterization versus
qualification. A temperature sweep asks how a healthy product behaves while it
is hot. A qualification stress asks whether the exposure caused permanent
change. Reversible movement and lasting degradation need different evidence and
support different decisions.

The second distinction that gets missed is
HTOL\sidenote{\textbf{HTOL}: high-temperature operating life. Representative
samples are biased at elevated temperature to accelerate mechanisms tied to
powered operation, then measured against a pre-stress baseline.} versus
burn-in\sidenote{\textbf{Burn-in}: an optional production screen that runs
every unit, or a sampled subset, under stress to cull a demonstrated
early-life population before shipment.}. Neither substitutes for the other.
HTOL is representative-sample life or mechanism evidence, so it is not an
every-unit ship screen. Burn-in is only justified when a specific early-life
mechanism is demonstrated and the screen is shown to remove it, because a
burn-in that consumes useful life while catching nothing is a net loss.

Reliability qualification is Step~6 of the product-readiness lifecycle in
\Cref{ch:product-readiness,sec:product-readiness}. Manufacturing validation and
ATP live in \Cref{ch:manufacturing}.

\section{Building the qualification argument}
\label{sec:qual-argument}

Qualification is a life and variation argument, not a list of stresses. Build
it in this order, and write it down in this order, because each step constrains
the next.

\begin{enumerate}
\item \textbf{State the claim.} Name the life, environment, handling exposure,
      and performance the product must hold, with the use condition and thermal
      class it applies to.
\item \textbf{Identify credible mechanisms.} Ask what physical or electrical
      process could violate that claim, at the die, the package, the assembly,
      the passive optics, and the connector.
\item \textbf{Select an acceleration method per mechanism.} Choose a stress
      that speeds up that mechanism without creating unrelated failure physics
      (\Cref{sec:accel-validity}).
\item \textbf{Define the observable before the stress.} Decide which measurable
      change reveals degradation, and at which reference plane, while you can
      still take a trusted baseline.
\item \textbf{Define acceptance before the stress.} State how much change is
      allowable and why that limit still closes the claim.
\item \textbf{Choose representative samples.} Cover lots, date codes,
      suppliers, sites, process corners, and product variants
      (\Cref{sec:qual-sample-strategy}).
\item \textbf{Run with intermediate reads.} A unit can stay functional while
      steadily consuming margin, and end-point-only data cannot show a trend.
\item \textbf{Interpret with stated confidence.} Connect observed drift, sample
      size, censoring, and model assumptions to a bounded statement.
\item \textbf{Decide, then hand off.} Release, restrict, derate, redesign, or
      gather more evidence, and name what production must control afterward
      (\Cref{ch:manufacturing}).
\end{enumerate}

Debugging is what you do when remaining margin hits zero. Qualification asks
how much margin remains after the expected stresses.

\engcheck{Margin budgeting}{%
Every stress spends margin: temperature, voltage, ripple, contamination,
insertion loss, connector wear, vibration, aging, process variation.
Qualification verifies what remains, not only that the part still links.}

The decision tree in \Cref{sec:tree-qual-evidence} walks the same sequence as a
set of exit criteria, and \Cref{fw:qual-plan} is the spoken interview form.

\section{Accelerated testing and model validity}
\label{sec:accel-validity}

Acceleration is the only reason qualification fits inside a program schedule,
and it is where most qualification arguments break. The requirement is
mechanism preservation: the stress must speed up the same physics that will act
in the field, and nothing else.

A severe test is not automatically a predictive test. Push temperature high
enough and you stop accelerating active-region wear and start melting solder or
degrading an epoxy that would never move at use conditions. That is mechanism
substitution, and the resulting number projects nothing. The same trap appears
at the other end: an overstress screen that kills units tells you the part has
a limit, not that the limit is reached in service.

Separate reversible movement from permanent change. Wavelength shifting with
temperature and recovering on cool-down is behavior. Threshold current that
comes back higher than it started is degradation. Only the second one belongs
in a life projection.

Arrhenius\sidenote{\textbf{Arrhenius}: acceleration factor
$\exp[(E_a/k)(1/T_\mathrm{use} - 1/T_\mathrm{stress})]$ for a single
thermally activated mechanism. The activation energy $E_a$ belongs to that
mechanism, not to the product.} projection applies when the named mechanism is
temperature-activated in the assumed regime. Take $E_a$ from mechanism-specific
supplier data, credible literature, or degradation data fitted to the actual
process, and report its uncertainty. One $E_a$ does not cover laser wear,
corrosion, solder fatigue, and electrical overstress at once. The physics of
laser aging behind the laser-side numbers is developed in
\Cref{sec:laser-aging,ch:lasers}.

Use the right temperature. Chamber air is not junction or active-region
temperature, and the difference between them is the difference between a
defensible acceleration factor and a wrong one by a large multiple. State which
temperature the model uses and how it was determined.

Two other named methods commonly appear alongside HTOL, and each answers a
narrower question. HTSL\sidenote{\textbf{HTSL}: high-temperature storage life.
An unbiased bake that targets thermal and material change without powered
operation, so it cannot substitute for HTOL.} isolates unbiased thermal and
material effects. HAST\sidenote{\textbf{HAST}: highly accelerated stress test.
Temperature, humidity, and pressure combined to accelerate moisture ingress in
far less time than damp heat, at the cost of a harsher validity argument.}
compresses moisture exposure. Method-by-method detail is in
\Cref{tab:stress-method-reference}.

Finally, plan the read points. Intermediate reads convert a pass/fail into a
degradation trend, and a trend supports projection while an endpoint does not.

\section{Samples, confidence, and evidence sufficiency}
\label{sec:qual-sample-strategy}

Weak answer: ``We test 20 units.'' Strong answer: the sample plan follows from
the failure-rate or degradation target, the required confidence, the cost of
the units, and the variation in the population you intend to ship. State lots,
date codes, suppliers, manufacturing sites or lines, process corners, and
whether the claim is a zero-failure upper bound or an observed rate.

\keyidea{Zero failures do not prove a zero failure rate.
0 failures in 20 units is not the same evidence as 0 failures in 1{,}000 units.
0 failures after $10^{3}$ device-hours is not the same evidence as 0 failures
after $10^{6}$ device-hours. Evidence strength depends on sample-hours,
one-sided upper confidence bounds, censoring, and whether lots and sites are
representative. State the bound, or do not make the claim.}

Sample-hours, not unit count, is the currency. Units multiplied by exposure
hours, adjusted by a justified acceleration factor, sets how tight the upper
bound can be. Unit count separately sets how much population variation you saw,
which is why one lot for a long time and many lots for a short time answer
different questions. When evidence is thin, a restricted release with a stated
residual risk is a legitimate outcome: ship a bounded population, keep the
monitoring in place, and widen only when the exposure supports it. Sizing rules
and censoring mechanics are in \Cref{sec:qual-sample-reference}.

\tradeoff{More qualification vs faster release}{Confidence and fewer late
escapes}{Schedule, cost, and delayed learning in the field}{When remaining risk
is high and reversible controls cannot cover it}{Prioritize tests by risk
$\times$ uncertainty $\times$ impact. Do not test everything equally.}

\subsection{Reliability rate versus quality rate}
\label{sec:fit-dppm}

A fleet reliability claim and a manufacturing quality claim use different
denominators: operating exposure for FIT and units evaluated at a defined
boundary for DPPM. They answer different questions and should not be compared
or substituted for one another (\Cref{sec:tree-evidence-block}).

FIT\sidenote{\textbf{FIT}: failures in time. One FIT is one failure per
$10^{9}$ device-hours. Always name the population unit, failure definition,
observation period, exposure, confidence, and whether the value is observed or
predicted. FIT commonly describes an approximately constant-hazard useful-life
regime; it is not automatically a complete lifetime model.} is an
exposure-normalized failure rate. One FIT means one failure per $10^{9}$
device-hours. A FIT value is meaningful only after the population unit is
named (for example, laser die, optical module, lane, or link) and the exposure
window, failure definition, censoring rules, confidence interval, and hazard
assumption are stated. The analysis must also distinguish an observed
fleet rate from a predicted rate derived from qualification data or a
reliability model.

A statement such as ``the module is 50~FIT'' is therefore incomplete. It should
identify what counted as a failure, how many eligible units were observed, for
how many device-hours, whether repaired or replaced units re-entered the
population, and whether the estimate assumes approximately constant hazard
during useful life. FIT should not be extended automatically into early-life or
wear-out regions where the hazard rate may change.

DPPM\sidenote{\textbf{DPPM}: defective parts per million. A unit-based quality
measure evaluated at a defined inspection, shipment, customer, or field
boundary. Its meaning depends on the denominator, time window, inspection
coverage, sampling, and treatment of retest, rework, and repeated returns.} is
a unit-normalized quality rate. It reports the number of defective units per
million units evaluated at a defined boundary. That boundary might be incoming
inspection, final test, outgoing quality inspection, customer receipt, or
confirmed field escape. The denominator, inspection coverage, time window,
sampling method, and treatment of rework, retest, and duplicate returns must be
stated.

A factory can therefore have low outgoing DPPM while the fleet has poor
long-term reliability, or high incoming DPPM while strong screening prevents
customer escapes. Conversely, a low FIT estimate does not prove that production
is free of assembly defects or process excursions.

The practical distinction is:
\begin{itemize}
\item FIT asks how failure occurrence scales with accumulated operating
      exposure.
\item DPPM asks what fraction of units is defective at a named quality
      boundary.
\end{itemize}

This chapter owns the reliability claim, exposure model, and confidence limits.
\Cref{ch:manufacturing} owns detailed DPPM accounting, yield decomposition,
screening effectiveness, and escape analysis.

\subsection{The three failure regimes}
\label{sec:bathtub}

Field failures arrive on three different clocks. Treat the regimes as
hypotheses that a failure pattern can discriminate between, not as confirmed
mechanisms.

\begin{description}
\item[Early life] A cluster concentrated in new units or a single date code
      suggests latent manufacturing defects. That points at screening,
      first-article control, and process containment. Burn-in is one possible
      response, not an automatic fix, and it is only justified when the
      early-life population is demonstrated.
\item[Useful life] A steady trickle with no date-code or age correlation
      suggests roughly constant-hazard behavior. That points at redundancy,
      fast repair, and spares rather than at a supplier corrective action.
\item[Wear-out] A rate that rises with install age suggests physics-driven
      degradation. That points at derating, life projection, and replacement
      planning. Connector wear and contamination accumulate with mate cycles
      and are usage-driven, so budget them separately from a flat-hazard rate
      (\Cref{sec:connector-reliability}).
\end{description}

The regime is an inference from the failure distribution, and it must be
confirmed against a mechanism before corrective action starts
(\Cref{ch:failure-modes}). The system consequence of each regime differs too: a
steady random rate is absorbed by a resilient fabric, while a correlated
early-life cluster is not (\Cref{sec:fabric-availability,ch:networking}).

\section{Mechanism families and qualification evidence}
\label{sec:wearout-modes}

Organize a qualification plan by mechanism family, not by test list and not by
field triage bucket. Each family below has its own physics, its own credible
acceleration, its own observable, and its own acceptance basis. Field
classification of a failure that already happened is a different job and lives
in \Cref{sec:fleet-triage}.

\begin{table*}[ht]
\refstepcounter{table}\label{tab:qual-mechanism-families}\label{tab:wearout-map}%
\centering
\setlength{\tabcolsep}{3pt}%
\footnotesize
\renewcommand{\arraystretch}{1.25}%
\begin{tabular}{@{}p{0.17\linewidth}p{0.20\linewidth}p{0.19\linewidth}p{0.21\linewidth}p{0.19\linewidth}@{}}
\toprule
Mechanism family & Credible threat & Qualification exposure & Observable degradation & Acceptance basis \\
\midrule
Powered semiconductor degradation & Active-region wear, facet damage, modal or EAM shift & Biased elevated temperature (HTOL) & $I_\mathrm{th}$, slope, $\lambda$, SMSR, TDECQ, bias creep & Projected drift still closes the life claim \\
\addlinespace[1.2em]
Thermo-mechanical fatigue & Expansion mismatch in solder, bonds, epoxy, attach & Temperature cycling; shock and vibration & Continuity, intermittent lanes, coupling or alignment step & No permanent step; no intermittents under monitoring \\
\addlinespace[1.2em]
Moisture, corrosion, material degradation & Ingress, electrochemical attack, seal or adhesive loss & Damp heat or HAST, biased when electrochemical & Leakage, insertion loss, ORL, physical evidence & Bounded loss and leakage; no seal breach \\
\addlinespace[1.2em]
Mechanical and optical-interface durability & Mate-cycle wear, debris, ferrule damage, handling & Mate cycling; shock; handling sequence & Insertion loss growth, ORL, endface grade & Loss and ORL inside the plant budget at rated cycles \\
\addlinespace[1.2em]
Electrical overstress, ESD, latch-up & Handling discharge, overvoltage, current injection & HBM/CDM classification; JESD78 injection & Classification level; supply-current latch; hard fail & Rated level covers the handling and system environment \\
\bottomrule
\end{tabular}
\par\medskip
{\raggedright\footnotesize\textbf{Table~\thetable.} Mechanism families and the
evidence each one needs. Plan coverage against this table before choosing
stresses. Named methods: \Cref{tab:stress-method-reference}. Filled planning
cells including production control: \Cref{tab:qual-planning-matrix}. Field
classification of an observed failure: \Cref{sec:fleet-triage}.\par}
\end{table*}
\FloatBarrier

\subsection{Powered semiconductor degradation}

Biased high-temperature exposure targets gradual active-region wear, facet
damage, SMSR collapse, and EAM absorption-curve shift. Observables that move
before hard failure include threshold, slope, wavelength, SMSR, and EAM bias
creep. On a bookended module, use launch power, OMA, wavelength, BER, supply
current, and telemetry, and narrow the claim to what those proxies support.
Acceptance is projected drift that still closes the life and system-margin
claim, not a final functional pass. Device physics and Arrhenius detail:
\Cref{sec:laser-aging,ch:lasers}. Acceleration limits:
\Cref{sec:accel-validity}.

\subsection{Thermo-mechanical fatigue}

Expansion mismatch among solder, bonds, attach, epoxy, and fiber alignment
produces opens, intermittents, or coupling steps. Temperature cycling is the
primary exposure; shock and vibration cover handling and repeated excitation.
Fatigue opens are often transient at temperature extremes and disappear at room
temperature, so in-situ or periodic monitoring is required when intermittents
are credible. Acceptance is no permanent step and no monitored intermittents.
Fixes are usually process controls in \Cref{ch:manufacturing}.

\subsection{Moisture, corrosion, and material degradation}

Damp heat and HAST target ingress, electrochemical attack, contamination
movement, and seal or adhesive degradation, not waterproofing. Bias when the
mechanism is electrochemical. Observables include leakage, insertion loss, ORL,
BER or functional change, and physical evidence. Acceptance is bounded loss and
leakage with no seal breach against the claimed storage and operating
environments. If HAST creates a field-irrelevant mechanism, the result is not
projectable.

\subsection{Mechanical and optical-interface durability}

Connector mate cycling, shock, and handling sequences address wear, debris,
alignment repeatability, insertion-loss growth, and ORL growth. Reflections
often move BER before average power moves. Acceptance uses plant loss and ORL
budgets at a service-based cycle count, not a generic datasheet rating. Method
detail: \Cref{sec:connector-reliability,app:reliability-reference}.

\paragraph{Photonic packaging and attach.}
\label{sec:photonic-packaging}
After lasers are screened, packaging and attach often dominate returns: FAU
alignment, epoxy creep, solder voids, underfill cracks, and TEC or thermal-path
failures that look like wavelength unlock
(\Cref{sec:cpo-status,sec:tfln-mzm,ch:wdm}).

\subsection{Electrical overstress, ESD, and latch-up}

Driver, TIA, retimer, and DSP damage from handling discharge or latch-up is not
an Arrhenius wear-out mechanism. Treat both as required electronics
qualification checkpoints, not as competitors with laser aging,
thermo-mechanical fatigue, moisture, or optical-interface durability for the
life argument.
HBM/CDM\sidenote{\textbf{HBM/CDM}: human-body model and charged-device model ESD
classification methods that set handling limits.} ratings classify handling
ESD. A damaged driver can look like a dead laser: no light, no optical aging
trend, supply-current anomaly, often a date-code cluster.
DPA\sidenote{\textbf{DPA}: destructive physical analysis used to confirm a
mechanism after external evidence is exhausted.} separates facet damage from
electrical damage when needed. Latent ESD rarely has an economical every-unit
screen, so handling controls and supplier evidence carry that part of the
argument (\Cref{sec:ic-reliability}).

Latch-up qualification checks whether an IC enters a parasitic high-current
state when its supply or I/O pins experience specified current or overvoltage
stress. It is relevant to module electronics, especially during hot-plug, power
sequencing, and interface transients. Latch-up is an electrical-susceptibility
mechanism rather than an Arrhenius wear-out process. Passing a component-level
latch-up method supports IC robustness, but it does not establish complete
module immunity to system-level sequencing, rail disturbance, or connector
events.

\begin{table*}[ht]
\refstepcounter{table}\label{tab:latchup-checkpoint}%
\centering
\setlength{\tabcolsep}{3pt}%
\footnotesize
\renewcommand{\arraystretch}{1.2}%
\begin{tabular}{@{}p{0.24\linewidth}p{0.36\linewidth}p{0.34\linewidth}@{}}
\toprule
Method & Contribution & Limitation \\
\midrule
JESD78-style latch-up testing & Evidence that an IC tolerates defined supply
and I/O injection stresses without entering a destructive high-current state &
Does not reproduce every module-level hot-plug, sequencing, grounding, or
transient condition \\
\bottomrule
\end{tabular}
\par\medskip
{\raggedright\footnotesize\textbf{Table~\thetable.} Latch-up as an electronics
qualification checkpoint. Test classes, injection procedures, temperature
conditions, and acceptance limits:
\Cref{sec:ic-reliability,app:reliability-reference}.\par}
\end{table*}
\FloatBarrier

\section{One worked qualification argument}
\label{sec:qual-worked-laser}

Walk one mechanism end to end before opening the matrix.

\textbf{Claim.} Five years of operation at the use condition and thermal class
named in the PRD, with the launch power and extinction the link budget assumes.

\textbf{Mechanism.} Gradual active-region wear: threshold current rises and
slope efficiency falls over life. This is not catastrophic optical damage,
which appears as sudden dark after a healthy ship LIV, and it is not ESD or
latch-up, which produce a hard fail with a supply-current signature and no
optical aging trend.

\textbf{Stress.} HTOL at an elevated temperature and bias chosen because they
accelerate that specific wear mechanism, with the junction temperature
determined rather than assumed from chamber air.

\textbf{Observable.} LIV against the ship baseline at named planes: threshold
current, slope efficiency, wavelength, and either launch power or BER drift on
bookended units. Intermediate reads at planned intervals so the result is a
trend, not an endpoint.

\textbf{Acceptance.} A bounded drift limit that, after projection to the use
condition, still closes the life claim with system margin intact. The
activation energy and its confidence are documented with the projection.

\textbf{Samples and confidence.} Representative lots, date codes, and sites,
with sample-hours sized to the target rate and confidence, and censoring stated
(\Cref{sec:qual-sample-strategy}).

\textbf{Decision and handoff.} Release, restrict, or derate based on the
projected drift and residual risk, then name the production control:
a room-temperature LIV or power proxy in the ATP, a sampled hot audit, SPC on
threshold and slope, or an explicit statement that no cost-effective 100\%
screen separates the weak tail. A production burn-in may cull infant mortality;
it does not replace this life argument.

Solder fatigue, contamination, and connector wear follow the same template with
different physics. The filled-in matrix for all families is
\Cref{tab:qual-planning-matrix} in \Cref{app:reliability-reference}. It is the
reference form of this story, not a substitute for telling it.

\section{Standards as evidence sources}
\label{sec:qual-standards}

Standards give you established stress methods, shared vocabulary, and supplier
reports you can accept without renegotiating a plan on every design win. They
do not establish that the product is suitable, that the relevant mechanisms
were covered, that the samples represented the shipping population, or that the
life claim closes. That argument is yours.

\begin{table*}[ht]
\refstepcounter{table}\label{tab:qual-standards-ownership}%
\centering
\setlength{\tabcolsep}{3pt}%
\footnotesize
\renewcommand{\arraystretch}{1.25}%
\begin{tabular}{@{}p{0.26\linewidth}p{0.34\linewidth}p{0.34\linewidth}@{}}
\toprule
Standard family & Owns & Scores on \\
\midrule
GR-468 / GR-3013 & Active optoelectronics: laser die, photodiode, modulator,
and their package (GR-3013 for shorter-life information-handling claims) &
LIV, spectral, and functional limits \\
\addlinespace[1.0em]
GR-1221 / GR-1209 & Passive optics: connectors, couplers, filters, MUX/DEMUX,
isolators (GR-1209 for functional and network-use criteria) & Insertion loss
and return loss \\
\addlinespace[1.0em]
JEDEC JESD47 & Driver, TIA, retimer, and DSP silicon qualification flow &
Parametric and functional limits after stress \\
\addlinespace[1.0em]
HBM/CDM & ESD classification for ICs & Rated level by pin and model \\
\addlinespace[1.0em]
JESD78 latch-up & IC supply and I/O injection robustness
(\Cref{tab:latchup-checkpoint}) & Module hot-plug, sequencing, and rail
immunity \\
\addlinespace[1.0em]
IEC connector methods & Mate-cycle durability and endface grading & Loss,
reflection, and endface pass zones \\
\bottomrule
\end{tabular}
\par\medskip
{\raggedright\footnotesize\textbf{Table~\thetable.} Which document owns which
part of the link. Split the qualification the way you split the failure budget.
Expanded ownership table and method detail:
\Cref{tab:gr-standards-ownership,sec:gr468,sec:ic-reliability,sec:latch-up,app:reliability-reference}.\par}
\end{table*}
\FloatBarrier

If a standard contains a test you believe is low risk, either run it or
document a waiver with mechanism reasoning, prior evidence, customer
agreement, and a named risk owner. Silently deleting the row is how a
qualification plan loses its meaning.\sidenote{AEC-Q100, the automotive IC
qualification standard, is sometimes raised in datacenter reviews. A published
Q100 grade is useful supplemental evidence about ESD and temperature-cycle
margin, but it is not a datacenter requirement and its absence is not a
finding (\Cref{sec:ic-reliability}).}

\section{Decisions and handoffs}
\label{sec:qual-handoffs}

Qualification exists to unlock a decision. Name which one, and state the
residual risk that comes with it.

\begin{itemize}
\item Release for unrestricted deployment.
\item Restricted release: a bounded population, site, or use condition.
\item Derate: narrow the operating envelope so the mechanism stays slow.
\item Redesign: the mechanism cannot be controlled at the required margin.
\item Supplier restriction: qualify a subset of lots, sites, or part numbers.
\item Additional qualification: extend exposure, add samples, or add read
      points before deciding.
\item Production control: a screen, ATP limit, or SPC chart that catches the
      mechanism at scale.
\item Sampled audit: periodic stressed sampling when a 100\% screen is not
      economical.
\item Fleet monitoring: a telemetry signature and a threshold that trips
      review.
\item No claim: state plainly that the evidence does not support the life
      statement.
\end{itemize}

Qualification identifies which mechanisms need control. It does not choose the
control. \Cref{ch:manufacturing} decides which production screen, supplier
requirement, or SPC chart is economical at volume, and whether a
room-temperature proxy correlates well enough to stand in for a stressed
measurement.

A component FIT is not a system availability number. Converting one into the
other requires redundancy, detection and reroute behavior, repair time, spares,
and the workload cost of each failure
(\Cref{sec:fabric-availability,ch:networking}).

When a qualification test fails, preserve the stressed state, the baselines,
the read history, the configuration, and the sample genealogy before anyone
retests. Scope whether the failure is isolated, lot-correlated,
stress-dependent, or systematic, then confirm the mechanism before choosing a
response (\Cref{ch:failure-modes}).

\paragraph{From qualification to production and fleet evidence.}
Qualification bounds life risk for representative hardware under stated
assumptions. Three things then happen in sequence. Production has to reproduce
that population, which is a capability and control problem, not a repeat of the
qualification (\Cref{ch:manufacturing}). The fleet accumulates the exposure no
qualification can afford, so field data is the first real test of the
mechanism model, with install age, temperature, firmware, supplier lot, and
return code recorded so regimes do not get mixed. Failures that pass
qualification and fail in the field usually sit in derating policy, connector
contamination, or a manufacturing coverage gap
(\Cref{sec:laser-aging,sec:fleet-triage,sec:escaped-defects}). Each of those
sends a specific correction back into the qualification plan for the next
product, which is how the argument improves rather than repeats.

\section{Interview takeaway}
\label{sec:reliability-design-review}

\keyidea{Reliability qualification is a mechanism-driven confidence argument.
Begin with the lifetime or environmental claim, identify how it could be
violated, select an acceleration method that preserves the mechanism, define
observable degradation and acceptance before the stress, and choose
representative samples. The result supports a bounded decision with stated
residual risk; it does not guarantee every production unit. Production controls
and fleet evidence carry the argument forward
(\Cref{ch:manufacturing,ch:failure-modes}).}

Junior mistake: treat zero fails in a small HTOL lot as a FIT claim, treat a
temperature sweep as temperature qualification, or read a completed standards
list as a life argument.

Better practice: state the claim first, name one mechanism per stress, declare
the observable and acceptance limit before the chamber door closes, and quote
the upper bound with its confidence rather than the number of units that passed
(\Cref{sec:qual-sample-strategy,sec:qual-terms}).

\subsection{Interview Q\&A: Reliability Qualification}
\label{sec:reliability-interview-qa}

Practice speaking these answers aloud. Prefer first-person reasoning over
stress inventories. Detail lives in
\Cref{sec:qual-argument,sec:qual-terms,sec:fit-dppm,tab:qual-planning-matrix,app:reliability-reference}.

\paragraph{Question 1. What is reliability qualification, and how does it
differ from system validation, manufacturing validation, and production
screening?}
\textit{Tests:} terminology, evidence ownership, and decision boundaries.

\textit{Spoken answer.}
``I treat reliability qualification as a bounded confidence argument that a
defined design will continue meeting named requirements through its claimed
life and environmental exposure. I build that argument from credible
degradation mechanisms, relevant acceleration, observable permanent change,
representative samples, and explicit model limits. System validation answers a
different question: whether the complete product is suitable for its intended
use across the supported system envelope. Manufacturing validation asks whether
production can repeatedly build, measure, trace, and control the qualified
design. Production screening asks whether a particular unit or production
population shows a detectable defect. The same measurement can support more
than one activity, but the claim, population, exposure history, acceptance
criterion, and resulting decision are different''
(\Cref{ch:product-readiness,ch:manufacturing,sec:qual-terms}).

\textit{Pressure follow-up.} ``Can the same temperature test support both
characterization and qualification?''\\
\textit{Answer pivot.} ``Yes, provided I separate reversible operating behavior
from permanent exposure-driven change. A temperature sweep of a healthy unit
characterizes how power, wavelength, BER, or control demand move while the unit
is hot. An extended biased exposure can support qualification when it targets a
named degradation mechanism and I evaluate lasting change after the device
returns to a defined measurement condition.''

\textit{Pressure follow-up.} ``Why isn't a long temperature sweep automatically
qualification?''\\
\textit{Answer pivot.} ``Because duration and severity are not enough. I still
need a life or environmental claim, a credible mechanism, a justified
acceleration model, representative samples, predefined observables, and an
acceptance rule tied to the product margin.''

\textit{Trap:} ``Qualification is system validation performed for longer, at
higher temperature, and with harsher tests.''

\paragraph{Question 2. How do you build a mechanism-driven qualification plan?}
\textit{Tests:} planning from product claims and credible risks rather than from
a standards checklist.

\textit{Spoken answer.}
``I start with the claims the product must support: operating life, storage,
environmental exposure, handling, and retained performance. Then I identify the
credible mechanisms that could violate those claims across the laser,
electronics, package, passive optics, and connectors. For each mechanism, I
choose an exposure that accelerates or reproduces it without substituting
unrelated failure physics. Before testing, I define the observable degradation,
measurement condition, acceptance criterion, sample population, and model
assumptions. I include representative lots, suppliers, assembly sites, and
process corners rather than testing many units from one convenient lot. I also
define the decision for each possible outcome: release, restricted release,
derating, redesign, supplier action, or additional evidence. Standards provide
established methods and a common language, but the mechanism-to-claim argument
determines why each test belongs''
(\Cref{sec:qual-argument,tab:qual-planning-matrix}).

\textit{Pressure follow-up.} ``The standard contains a test you believe addresses
a low-risk mechanism. Do you omit it?''\\
\textit{Answer pivot.} ``I would not silently remove it. I first check whether
it is contractually required, whether the mechanism is genuinely absent, and
whether prior or supplier evidence applies to the released design and process.
Then I either run the test or document a formal waiver that states the technical
basis, evidence, residual risk, and accountable approver.''

\textit{Pressure follow-up.} ``What makes an accelerated stress invalid?''\\
\textit{Answer pivot.} ``It is invalid for the claim when it creates a different
dominant mechanism from the field condition, when the acceleration model is
unsupported, or when the measured observable does not reveal the degradation
that threatens the requirement. A harsher test is not automatically a more
predictive test.''

\textit{Trap:} ``I select the applicable standard, run every listed stress, and
declare the product qualified when all rows pass.''

\paragraph{Question 3. What is the difference between a failure mode and a
failure mechanism, and why does failure timing matter?}
\textit{Tests:} distinguishing the observed symptom from the underlying physics,
then using timing to identify the failure regime.

\textit{Spoken answer.}
``A failure mode is what I observe: for example, low optical power, high BER,
intermittent operation, or a dead lane. A failure mechanism is the physical or
electrical process that caused the symptom, such as active-region degradation,
optical-coupling shift, corrosion, solder fatigue, contamination, or ESD damage.
One mode can result from several mechanisms, so the symptom alone does not tell
me what to qualify or correct.

Failure timing provides a second discriminator. A cluster of early failures
tied to a particular date code points toward a manufacturing escape. A roughly
steady failure rate with little dependence on age is consistent with
constant-hazard, random failure behavior. A rate that increases with accumulated
age or exposure points toward wear-out. Timing does not prove the mechanism, but
it narrows the hypotheses and helps determine whether the appropriate response
is additional qualification, process CAPA and containment, failure analysis, or
service-life and replacement planning''
(\Cref{tab:qual-mechanism-families,sec:bathtub,ch:failure-modes}).

\textit{Pressure follow-up.} ``Give me two mechanisms that can cause low optical
power.''\\
\textit{Answer pivot.} ``Laser degradation can increase threshold current or
reduce slope efficiency, causing the source itself to generate less power.
Alternatively, coupling movement, connector contamination, or increased path
loss can reduce delivered power even when the laser remains healthy. I would
distinguish them using source-level measurements, electrical-drive and control
telemetry, named-plane optical power, insertion loss, and physical evidence. The
symptom is the same, but the evidence and corrective controls are different.''

\textit{Pressure follow-up.} ``A new date code has many early failures. Is that
proof of infant mortality?''\\
\textit{Answer pivot.} ``No. It is a strong population clue, not a confirmed
mechanism. I would first verify exposure, configuration, denominator, and
failure definition, then compare genealogy and physical evidence before
assigning a process cause or adding a screen.''

\textit{Trap:} ``Low optical power means the laser is aging.''

\paragraph{Question 4. When can accelerated testing support a field-life claim?}
\textit{Tests:} mechanism-specific acceleration, model validity, and
uncertainty.

\textit{Spoken answer.}
``Accelerated testing can support a field-life claim only when the applied
stress accelerates the same dominant failure mechanism expected under field
conditions, without introducing a different mechanism or changing the
rate-limiting physics.

For a thermally activated mechanism, an Arrhenius model may translate stress
exposure into equivalent use-condition exposure. However, the activation
energy, temperature range, and other model parameters must be appropriate to
that specific mechanism and process. I use the temperature at the relevant
physical location (such as junction or active-region temperature) rather than
chamber air alone. I also define the degradation signature and failure criterion
in advance, confirm that stressed and field-relevant samples exhibit consistent
physical evidence, and account for duty cycle, operating conditions, sample
variability, and statistical uncertainty.

The result is a bounded life estimate under stated assumptions, not a universal
conversion from test hours to field years. A more severe test is not
automatically a more predictive test''
(\Cref{sec:accel-validity,sec:laser-aging}).

\textit{Pressure follow-up.} ``How do you choose the activation energy?''\\
\textit{Answer pivot.} ``I choose it from evidence specific to the mechanism
and, ideally, to the actual materials and process: supplier characterization,
credible published data, or degradation measurements collected at multiple
stress levels. I also report the sensitivity of the life estimate to uncertainty
in that value. I would not reuse one activation energy across unrelated
mechanisms such as laser degradation, corrosion, and solder fatigue.''

\textit{Trap:} ``Hotter means faster, so I can convert test hours directly into
field years.''

\paragraph{Question 5. Walk me through HTOL for a laser-bearing optical
module.}
\textit{Tests:} mechanism-driven powered aging, access-aware observables, and
margin-based acceptance.

\textit{Spoken answer.}
``HTOL applies electrical bias at elevated temperature to accelerate credible
mechanisms associated with powered operation. I begin with the field-life claim
and the mechanism I intend to accelerate, then choose the stress temperature,
bias, duty cycle, duration, and sample population. The operating point must be
severe enough to provide useful acceleration without creating failure physics
that would not dominate in service. I establish a stable baseline at defined
measurement conditions and determine junction or active-region temperature
rather than treating chamber air as the device temperature.

With engineering access, I track LIV behavior, threshold current, slope
efficiency, wavelength, SMSR, and selected RIN measurements. On a bookended
module, I may be limited to launch power, OMA, wavelength, BER, supply current,
monitor and control telemetry, and remaining control-loop headroom. In that
case, I correlate those external observables with engineering-access or supplier
data and limit the claim to what they can actually reveal.

I include intermediate read points and distinguish measurements made under
stress from standardized post-stress measurements, so reversible temperature
effects are not mistaken for permanent degradation. Acceptance is based on
predefined catastrophic-failure and parameter-drift limits tied to end-of-life
margin. A module that still links at the end has not necessarily passed if it
has consumed the margin needed to support the life claim''
(\Cref{sec:qual-worked-laser}).

\textit{Pressure follow-up.} ``You have no internal LIV access. Is the test
useless?''\\
\textit{Answer pivot.} ``No. I use the strongest external proxies available,
such as named-plane optical power, OMA, wavelength, supply current, control
telemetry, BER, and control-loop headroom. I preserve correlation to
engineering-access or supplier data and narrow the claim accordingly. A
bookended BER result can demonstrate retained function, but it cannot by itself
identify laser degradation or quantify the remaining laser margin.''

\textit{Trap:} ``Run the module hot for 1,000 hours; if it still links, it
passes HTOL.''

\paragraph{Question 6. What distinct failure mechanisms do temperature
cycling, mechanical shock, and vibration reveal?}
\textit{Tests:} separating thermomechanical fatigue, discrete overload, and
repeated dynamic excitation.

\textit{Spoken answer.}
``These are durability stresses, not simply functional tests performed at
temperature or while the module is moving. Temperature cycling repeatedly
drives coefficient-of-expansion mismatch among solder joints, wire bonds,
adhesives, fiber attach, package materials, and optical alignment features. It
can reveal fatigue, cracking, delamination, intermittent contact, and
alignment drift. The relevant severity depends on the temperature range, ramp
rate, dwell time, cycle count, and temperature gradients within the assembly.

Mechanical shock applies a short-duration acceleration pulse representative of
a discrete handling, shipping, or installation event. It can reveal fracture,
permanent displacement, connector or fiber damage, and marginal retention.
Vibration applies repeated dynamic excitation over a defined spectrum. It can
excite resonances and reveal fretting, loosening, fatigue, intermittent
electrical contact, or optical-power modulation.

I choose the axes, mounting fixture, stress profile, and powered or unpowered
state from the credible field exposure and mechanism. When intermittent
behavior is credible, I monitor continuity, lane alarms, optical power, and BER
in situ, then repeat named-plane measurements under standardized conditions and
inspect for physical evidence. A final functional pass alone can miss a
momentary open or alignment excursion that disappears when the stress is
removed''
(\Cref{tab:stress-method-reference}).

\textit{Pressure follow-up.} ``Would you measure only before and after
cycling?''\\
\textit{Answer pivot.} ``Not when interconnect or alignment intermittency is
credible. I use periodic or in-situ monitoring, especially during temperature
transitions and at the hot and cold extremes, and correlate detected events
with the stress history. I still make baseline and post-stress measurements at
a common reference condition to separate reversible behavior from permanent
change.''

\textit{Trap:} ``Those tests simply confirm that the module operates while it
is hot, cold, or shaken.''

\paragraph{Question 7. What do damp heat and connector mate cycling qualify,
and what evidence do you collect?}
\textit{Tests:} moisture-driven degradation, bias dependence, and
optical-interface durability.

\textit{Spoken answer.}
``Damp heat exposes the product to controlled temperature and humidity to
accelerate moisture ingress or absorption and the resulting corrosion,
electrical leakage, material degradation, delamination, contamination movement,
or loss of interface integrity. The temperature, humidity, duration, and bias
state must match the mechanism being evaluated. I use electrical bias when an
electric field is needed to drive electrochemical migration or corrosion; an
unbiased exposure does not make the same claim. I track leakage and supply
current, optical and electrical parameter drift, visual evidence, and selected
destructive analysis, with measurements at a defined reference condition to
separate reversible moisture uptake from permanent damage.

Connector mate cycling qualifies the repeatability and durability of the
optical and mechanical interface through repeated connection events. It can
reveal endface wear, scratches, debris generation, ferrule damage, alignment
variation, intermittent contact, insertion-loss growth, and degraded optical
return loss. I control the mating counterpart, cycle procedure, cleaning
policy, and inspection criteria, then track insertion loss and ORL at a named
reference plane, along with optical power, BER, and endface condition.

Neither test is cleared by stable average optical power alone. A connector can
retain low insertion loss while increased reflection or intermittent alignment
creates a BER or noise floor, and moisture damage can first appear in leakage or
control demand before causing a large power change''
(\Cref{sec:connector-reliability}).

\textit{Pressure follow-up.} ``BER worsens after repeated connector mating, but
average power barely moves. What do you check?''\\
\textit{Answer pivot.} ``I inspect both endfaces for debris, scratches, and
ferrule damage, then measure insertion loss and ORL at the defined reference
plane. I test whether BER, eye quality, or RIN changes with the reflection
condition by cleaning, remating, or substituting a known-good interface. Stable
average power rules out a large added loss, but it does not rule out reflection
penalty or intermittent alignment.''

\textit{Trap:} ``Damp heat proves the package is waterproof, and mate cycling
only checks whether the latch still works.''

\paragraph{Question 8. How do GR-468, GR-1221/GR-1209, and JESD47 divide the
qualification problem?}
\textit{Tests:} standards applicability, subsystem ownership, and the gap
between component evidence and a product-level claim.

\textit{Spoken answer.}
``I divide the evidence by subsystem and failure physics. GR-468 is the
reliability-assurance framework for active optoelectronic devices such as
lasers, photodiodes, modulators, and their applicable package levels. It
provides established accelerated-aging, environmental, and mechanical methods,
with observables such as LIV, spectral behavior, and retained function.

GR-1221, used with GR-1209, covers passive optical components such as
connectors, couplers, splitters, filters, MUX/DEMUX devices, and isolators.
GR-1221 contributes reliability-assurance evidence, while GR-1209 supplies
functional, environmental, and network-use criteria. These parts are generally
scored on insertion loss, optical return loss, and physical durability rather
than on laser observables.

JESD47 and its referenced JEDEC methods provide the baseline qualification flow
for the driver, TIA, retimer, DSP, and other ICs. The evidence addresses the
released die, process, and package through stresses such as HTOL, temperature
cycling, moisture exposure, and mechanical tests, with parametric and
functional acceptance criteria.

I therefore split the qualification the way I split the failure budget, but I
do not assume that the component reports qualify the assembled module. I still
verify that each report applies to the released supplier, process, package,
site, and change state, then close product-level gaps such as optical alignment,
thermal interaction, control-loop behavior, connector interfaces, hot-plug
events, and combined-stress effects. The standards provide established methods
and evidence; the product claim still has to connect mechanism, stress,
observable, sample population, acceptance criterion, and decision''
(\Cref{sec:qual-standards,tab:gr-standards-ownership,app:reliability-reference}).

\textit{Pressure follow-up.} ``Every component supplier has a qualification
report. Is the module qualified?''\\
\textit{Answer pivot.} ``No. I first confirm that each report covers the exact
released configuration and intended use. I then identify mechanisms introduced
by assembly and subsystem interaction, because no component report proves
module-level optical alignment, thermal coupling, power sequencing,
interoperability, or field life.''

\textit{Pressure follow-up.} ``Where do ESD and latch-up fit?''\\
\textit{Answer pivot.} ``They are complementary IC-level robustness
checkpoints, not Arrhenius wear-out rows. HBM and CDM classify component ESD
susceptibility, and JESD78 evaluates latch-up under defined pin-injection and
supply-overvoltage conditions. I translate those results into handling,
assembly, and process controls, but I still evaluate module-level system ESD,
hot-plug, rail disturbance, and power sequencing. Latent ESD often has no
reliable 100\% screen, so prevention and handling discipline matter''
(\Cref{tab:latchup-checkpoint,sec:latch-up}).

\textit{Trap:} ``The laser passed GR-468 and the electronics passed JESD47, so
the complete module is qualified.''

\paragraph{Question 9. Explain FIT, DPPM, sample size, and what zero failures
actually mean.}
\textit{Tests:} denominator discipline, exposure-based confidence bounds, and
population representation.

\textit{Spoken answer.}
``FIT and DPPM answer different questions because they use different
denominators. FIT is an exposure-normalized failure rate: one FIT is one failure
per $10^{9}$ device-hours. When I quote FIT, I name the population unit (die,
module, lane, or link), the failure definition, observation period, accumulated
exposure, censoring rules, confidence bound, hazard regime, and whether the
value is observed or model-predicted. FIT is commonly used for approximately
constant-hazard useful-life behavior; it is not automatically a complete
lifetime model.

DPPM is a unit-normalized quality measure at a defined boundary. I state whether
the count refers to incoming material, outgoing product, customer receipt, or
confirmed field escapes, and I define the denominator, inspection coverage,
time window, sampling, and treatment of rework, retest, and duplicate returns.
A product can have low outgoing DPPM and poor wear-out reliability, or high
incoming DPPM with effective screening and few customer escapes, so I do not
substitute DPPM for FIT.

Sample size follows from the claim. For a failure-rate claim, total equivalent
device-hours determine the statistical upper bound, but the acceleration factor
is usable only when its mechanism and model are justified. Unit count and the
distribution across lots, date codes, suppliers, sites, and process corners
separately determine how much population variation the test represents. For a
parameter-drift claim, I size the study from the expected variation, effect or
quantile of interest, confidence, and statistical power.

Zero failures therefore do not establish a zero rate. Under an independent,
constant-hazard Poisson model, zero failures over total equivalent exposure
$T$ give the one-sided upper bound
$\lambda_{\mathrm{U}}=-\ln(1-C)/T$ at confidence $C$. I report that bound and
its assumptions rather than saying that the tested units proved the field rate
is negligible''
(\Cref{sec:fit-dppm,sec:qual-sample-strategy}).

\textit{Pressure follow-up.} ``Twenty units each complete 1{,}000 device-hours
with zero failures. What can you claim at 90\% confidence?''\\
\textit{Answer pivot.} ``With no acceleration credit, the total exposure is
20{,}000 device-hours. Under the constant-hazard Poisson assumptions, the
one-sided 90\% upper bound is $2.303/20{,}000$, or about
$1.15\times10^{-4}$ failures per hour, approximately 115{,}000~FIT. A justified
acceleration factor could increase the equivalent use-condition exposure, but
twenty units would still provide limited coverage of population variation.''

\textit{Pressure follow-up.} ``How does component FIT become a fabric-level
risk?''\\
\textit{Answer pivot.} ``I first keep the population unit consistent. Under an
independent constant-hazard model, deployed population times FIT times exposure
gives the expected failure count. I then translate failures into service impact
using detection coverage, redundancy, reroute behavior, repair time, sparing,
and workload sensitivity, and I treat common-cause failures separately.
Component occurrence rate and system consequence are distinct parts of the
availability argument''
(\Cref{sec:fabric-availability,ch:networking}).

\textit{Trap:} ``Twenty units passed, so the field failure rate is effectively
zero and the lot is representative of production.''

\paragraph{Question 10. How do you define acceptance criteria, and when is
qualification evidence sufficient?}
\textit{Tests:} predeclared decision rules, life-margin closure, and
decision-scaled evidence.

\textit{Spoken answer.}
``I define acceptance before the stress from the requirement, the credible
mechanism, and the degradation budget available at end of life. For each test I
state the measurement condition and reference plane, baseline method, hard
failure and intermittency rules, allowable parameter shift or distribution,
physical-evidence criteria, and treatment of measurement uncertainty. Typical
criteria include no catastrophic failure, bounded drift in power or
sensitivity, no new BER or noise floor, preserved control-loop headroom, and no
mechanism-relevant physical damage.

Acceptance is not simply ``inside the data-sheet limit after stress.'' I compare
results at a common reference condition, account for guardband and measurement
uncertainty, examine individual trajectories and population shifts, and project
degradation to the claimed use condition only with a justified model. A unit can
remain functional at the endpoint yet fail the qualification argument because
it has consumed too much future margin.

Evidence is sufficient when the credible mechanisms are covered, stresses
preserve the intended physics, samples represent the released design and
process, the test execution is valid, and the exposure, trends, confidence
bounds, outliers, and model uncertainty support the decision being made.
Sufficiency is decision-dependent: a limited, monitored, reversible pilot can
carry more residual uncertainty than an unrestricted fleet release. I document
the exact claim supported, the claims not supported, the remaining risk, and
the product or process changes that would trigger requalification''
(\Cref{sec:qual-argument,sec:qual-sample-strategy,tab:qual-planning-matrix}).

\textit{Pressure follow-up.} ``All units pass the final limit, but one
parameter is trending steadily toward it. Pass or fail?''\\
\textit{Answer pivot.} ``The endpoint criterion passes, but the life claim is
still open. I first verify the measurement and examine the individual and
population trajectories. Then I test whether the trend is linear, saturating,
or accelerating and project it to the use condition with parameter and model
uncertainty. If the conservative projection crosses the end-of-life limit, I
fail or restrict the claim, extend the evidence, or correct the design. If it
closes with adequate margin, I can pass it with that margin and uncertainty
documented. I do not average away a credible degrading unit.''

\textit{Trap:} ``Every sample is still within its data-sheet limit after every
standard test, so qualification is complete.''

\paragraph{Question 11. What do you do when a qualification test fails?}
\textit{Tests:} test validity, as-found evidence, containment, mechanism
confirmation, and scoped requalification.

\textit{Spoken answer.}
``I first preserve the as-found evidence. I do not power-cycle, clean, rework,
or repeatedly probe the failed unit until I have captured its state, baseline
and intermediate reads, stress history, configuration and firmware, fixture
and instrument logs, control-unit results, and sample genealogy.

The first decision is whether I have a valid product failure or an invalid
test. I check the chamber profile, actual device stress, instrumentation,
calibration, fixture, controls, and protocol deviations. If the test was
invalid, I document why, correct the test-system cause, and repeat enough valid
exposure to support the original claim. I do not quietly discard the result or
call a retest pass corrective action.

If the failure is valid, I contain the potentially affected population and
freeze relevant release or change activity. I classify the observed mode,
timing, and population pattern, then confirm the mechanism using discriminating
measurements, controlled swaps, reproduction, and physical analysis in an
order that preserves evidence. The response may be a design or process change,
supplier CAPA, derating, restricted use, additional screening, or more
evidence.

Requalification follows a documented change-impact analysis. I repeat the
stresses needed to show that the mechanism and affected interfaces are
corrected, add regression coverage for risks introduced by the change, and
preserve applicable prior evidence. I do not automatically restart every
unrelated row, but I also do not rerun only the failed row when the corrective
action changes other failure physics''
(\Cref{sec:qual-argument,sec:accel-validity,ch:failure-modes}).

\textit{Pressure follow-up.} ``The failure occurred beyond the intended use
condition. Can you dismiss it?''\\
\textit{Answer pivot.} ``Not from stress level alone. I determine whether the
failure has the same physical signature as a field-relevant mechanism or
whether the excessive stress created a distinct mechanism that cannot occur
inside the qualified envelope. If it is the same mechanism, the result remains
relevant even though it appeared beyond the use condition. If it is demonstrably
an overstress-only mechanism, I can exclude it from the release claim with the
transition, evidence, and residual risk documented. If the boundary is
uncertain, the qualification argument remains open.''

\textit{Trap:} ``Rerun the unit; if it passes, record the first result as an
anomaly and continue qualification.''

\paragraph{Question 12. Give me a 60-second qualification plan for a new
optical transceiver.}
\textit{Tests:} concise synthesis, risk prioritization, evidence ownership, and
release handoff.

\textit{Spoken answer.}
``I first freeze the released configuration and its life, duty-cycle,
environmental, handling, and retained-performance claims. I map each claim to
credible mechanisms across the laser, electronics, package, passive optics,
connectors, and their interfaces, then rank the mechanisms by risk and
uncertainty. I reuse supplier GR-468, GR-1221, or JESD47 evidence only after
confirming that it applies to the released design, process, package, site, and
use condition.

For each remaining gap, I choose a stress that preserves the field mechanism
and predeclare the observable, reference measurement condition, intermediate
reads, catastrophic and drift limits, sample population, and confidence basis.
The plan typically combines powered aging, thermal and mechanical durability,
moisture and bias, connector cycling, and electronics robustness as required by
the architecture. Samples span representative lots, suppliers, sites, and
process corners.

I verify the applied stress, analyze trends as well as hard failures, and
confirm valid failures to a mechanism before corrective action. The release
gate is explicit: release, restrict, derate, redesign, or gather more evidence.
I hand the result to production as named screens, monitors, SPC variables,
supplier controls, and requalification triggers, and to the fleet as telemetry
and feedback. The output is a bounded supported claim with residual risk, not a
statement that the product is universally reliable''
(\Cref{sec:qual-argument,tab:qual-planning-matrix}).

\textit{Pressure follow-up.} ``Apply that to laser aging, and tell me whether
production burn-in replaces HTOL.''\\
\textit{Answer pivot.}
``For a five-year claim, I target active-region degradation with representative
qualification lots, justified biased high-temperature exposure, measured
active-region temperature, baseline and intermediate LIV or correlated external
proxies, and drift limits tied to end-of-life margin. I project only with a
mechanism-specific acceleration model and report its uncertainty.

Production burn-in serves a different regime. It can screen early-life
manufacturing defects only when the screen stress, duration, detection
signature, and escape correlation are validated. It does not provide the
sample-hours, population strategy, or degradation projection needed for the
wear-out claim, so it does not replace HTOL. I also verify that the screen does
not consume useful life or create damage that would not otherwise occur''
(\Cref{sec:qual-worked-laser}).

\textit{Trap:} ``I would collect the supplier reports, run the standard stress
list, and release the module if every row passes.''

Score each answer using the shared chapter-interview rubric in
\Cref{sec:chapter-interview-rubric}; repeat any answer that does not state a
requirement, mechanism, evidence, confidence, and decision.
